\section{Box-face integrals and analytic singularities}
\label{sec:face-exact}

For an unconstrained smooth objective, the root-box faces give a complete
description of certification complexity for a strictly positive separable
quadratic gap.  The description includes the full box.  It also includes
lower-dimensional faces, because a minimizer on the boundary can be sharp
in its inward normal direction and flat along the boundary.  A volume
integral over the full box can then miss the part of the domain that
requires the most subdivision.

Throughout this section, let
\[
 D=\prod_{i=1}^n[\ell_i,u_i],\qquad \ell_i<u_i,
 \qquad f_*=\min_D f,\qquad m=f-f_*.
\]
There are no constraints other than membership in $D$.  The function $f$
is continuously differentiable on a neighborhood of $D$, and its gradient
is Lipschitz on $D$.  Fix coefficients $\alpha_i>0$ such that
$f(x)+\sum_i\alpha_i x_i^2$ is convex on $D$.  For a test box
$B=\prod_i[a_i,b_i]\subseteq D$, use the exact separable $\alpha$BB bound
\begin{equation}
 \label{face:eq:oracle}
 L(B)=\min_{x\in B} f_B(x),\qquad
 f_B(x)=f(x)-\sum_{i=1}^n\alpha_i(x_i-a_i)(b_i-x_i).
\end{equation}
The assumed convexity makes $f_B$ convex.  Positive coefficients can
always be chosen: if the gradient has Lipschitz constant $M$, taking
$\alpha_i\ge M/2$ suffices, because the quadratic Taylor-remainder
bound gives a supporting plane for $f+\sum_i\alpha_i x_i^2$ at every
point of $D$.  The coefficients and the representation
\eqref{face:eq:oracle} remain fixed as the tolerance changes.
The separable quadratic underestimator is the established $\alpha$BB
construction \citep{adjimanDallwigFloudasNeumaier1998AlphaBBTheory}.

A box is valid at tolerance $\varepsilon>0$ when
$L(B)\ge f_*-\varepsilon$.  Here this has the exact pointwise form
\begin{equation}
 \label{face:eq:validity}
 m(x)+\varepsilon\ \ge\
 \sum_i\alpha_i(x_i-a_i)(b_i-x_i)
 \qquad(x\in B).
\end{equation}
We use the separate certificate quantities $N_{\mathrm{cover}}$,
$N_{\mathrm{rect}}$, and $N_{\mathrm{tree}}$: the minimum number of valid
boxes in a cover of $D$, in a rectangular partition of $D$, and among
the leaves of an axis-parallel split tree, respectively.  Thus
\begin{equation}
 \label{face:eq:benchmarks}
 N_{\mathrm{cover}}(\varepsilon)
 \le N_{\mathrm{rect}}(\varepsilon)
 \le N_{\mathrm{tree}}(\varepsilon).
\end{equation}
The upper construction uses the fixed optimal incumbent $f_*$.  Its cost
therefore measures certification after this incumbent is available.

Let $\mathcal F_+(D)$ denote all closed faces of $D$ of positive
dimension, including $D$ itself.  A face $F$ fixes some coordinates at
root endpoints; write $J_F$ for its free coordinates and
$d_F=|J_F|$.  Lebesgue measure in its affine hull is denoted by $\sigma_F$.
Define
\begin{equation}
 \label{face:eq:integral}
 I_F(\varepsilon)=
 \int_F(m(x)+\varepsilon)^{-d_F/2}\,d\sigma_F(x),\qquad
 \mathcal I_D(\varepsilon)=1+\sum_{F\in\mathcal F_+(D)}I_F(\varepsilon).
\end{equation}
All comparison constants below may depend on this fixed instance, its
dimension, and its coefficients, but not on $\varepsilon$.

\begin{theorem}[Box-face characterization]
 \label{face:thm:characterization}
Under the preceding assumptions, let $T_{\mathrm{dyad}}(\varepsilon)$
be the tree that halves every coordinate of each invalid box and tests
its $2^n$ children.  Then, as $\varepsilon\downarrow0$,
\begin{equation}
 \label{face:eq:characterization}
 N_{\mathrm{cover}}(\varepsilon)
 \asymp N_{\mathrm{rect}}(\varepsilon)
 \asymp N_{\mathrm{tree}}(\varepsilon)
 \asymp |T_{\mathrm{dyad}}(\varepsilon)|
 \asymp \mathcal I_D(\varepsilon).
\end{equation}
More explicitly, for every $\varepsilon>0$,
\begin{equation}
 \label{face:eq:lower-explicit}
 N_{\mathrm{cover}}(\varepsilon)
 \ge
 \max\left\{1,
 \max_{F\in\mathcal F_+(D)}
 \frac{d_F^{d_F/2}}{\pi^{d_F}}
 \left(\prod_{i\in J_F}\alpha_i\right)^{1/2}
 I_F(\varepsilon)\right\}.
\end{equation}
There are finite constants $C_0,C_1$, independent of $\varepsilon$, such
that
\begin{equation}
 \label{face:eq:upper-explicit}
 |T_{\mathrm{dyad}}(\varepsilon)|
 \le C_0+C_1\sum_{F\in\mathcal F_+(D)}I_F(\varepsilon).
\end{equation}
The same dyadic partition has a binary split-tree realization with at
most a constant multiple of this number of nodes.
\end{theorem}

The proof is in \cref{app:face:characterization}.  Its lower bound is an
arcsine integral on each root face.  For the upper bound, an invalid box
contains a point with $m+\varepsilon$ of order its squared side length.
Nonnegativity of $m$ on the closed root box controls the derivative
toward a nearby face, even when there is no room for an outward step.
Moving the point to that face gives a relative cube of comparable size
inside a slightly larger sublevel set.  Packing these relative cubes
and summing their scales gives the face integrals.  Optimal vertices
require a separate argument: either every inward edge derivative is
positive and the corner box is certified after a fixed refinement, or
an edge integral accounts for its remaining levels.  No separate
$\log(1/\varepsilon)$ term is needed for vertices.

The theorem assumes neither quadratic growth nor analyticity.  It also
does not identify the finite values of the three certificate minima.
The comparisons follow from a lower bound for arbitrary covers and an
upper bound constructed by a tree.  A total-node count and a leaf count
differ by a fixed factor for these finite trees; additional relaxation
evaluations or reduction rounds remain separate costs.

\begin{remark}[Gap assumptions and owned certificates]
 \label{face:rem:gap-scope}
The proof uses only two properties of the oracle: a bound implying
$m+\varepsilon\ge\alpha q_B$ on every valid box, with $\alpha>0$, and
an upper-error bound $f_B\ge f-\alpha' q_B$, with $\alpha'>0$ and
$q_B=\sum_i(x_i-a_i)(b_i-x_i)$.  Thus the characterization also holds
under these two bounds, with modified constants.  Strict positivity is
essential to this scope.  For example, choosing zero correction for a
convex objective gives an exact root bound and can give a one-box
certificate.

The per-face lower bound also applies to test boxes $B_e$ with measurable
owned sets $A_e\subseteq B_e$ covering $D$, when
\eqref{face:eq:validity} is known only on $A_e$.  The integral over each
owned part is bounded by the same full-box arcsine integral.  The event
accounting from the certificate framework can therefore transfer this
lower bound to allowed reductions.  Such a transfer bounds the specified
event or evaluation cost; it does not identify an owned family with a
valid closed-box cover.
\end{remark}

\subsection{Analytic objectives and exact power--logarithm rates}

Suppose now that $f$ is real analytic on a neighborhood of $D$.  On a
face $F\in\mathcal F_+(D)$ with a nonzero restriction
$h=m|_F$ and $\min_Fh=0$, define its real log canonical threshold and
multiplicity $(\lambda_F,\theta_F)$ through the first positive pole and
its order in
\begin{equation}
 \label{face:eq:zeta}
 \zeta_F(z)=\int_F h(x)^{-z}\,d\sigma_F(x).
\end{equation}
The integral initially converges for sufficiently small positive real
$z$, and the definition uses its meromorphic continuation.  The
restriction must be considered on $F$ as a domain in its own affine
hull.  A boundary zero is included, and no nonnegativity outside $F$ is
required.

The analytic input is classical.  On these compact semianalytic
domains, $\lambda_F$ is a positive rational number,
$\theta_F$ is a positive integer, and
\begin{equation}
 \label{face:eq:laplace}
 Z_F(s):=\int_F e^{-s h(x)}\,d\sigma_F(x)
 \sim c_{Z,F}s^{-\lambda_F}(\log s)^{\theta_F-1},
 \qquad c_{Z,F}>0.
\end{equation}
This is the compact-domain asymptotic of
\citet[Section~2, Theorem~2.10]{lin2017LearningCoefficient}.
His domain treatment resolves the analytic function and the boundary
inequalities together.  A closed face has affine analytic boundary
inequalities in its intrinsic coordinates and density one, so it
satisfies those hypotheses.  His absolute-value convention agrees with
\eqref{face:eq:zeta}, because $h\ge0$ on $F$.  A restriction that is
identically zero is handled separately below; it has no ordinary
positive-pole RLCT pair of this form.

\begin{lemma}[Evaluation of the regularized integral]
 \label{face:lem:analytic-integral}
Let $h\ge0$ be a nonzero analytic function on a positive-dimensional
compact box $F$, with a zero on $F$, and suppose its Laplace integral
satisfies \eqref{face:eq:laplace} with parameters
$\lambda,\theta,c_Z$.  For $a>0$, put
$I_a(\varepsilon)=\int_F(h+\varepsilon)^{-a}\,d\sigma_F$.  Then
\begin{align}
 I_a(\varepsilon)&\sim
 \frac{c_Z\Gamma(a-\lambda)}{\Gamma(a)}
 \varepsilon^{\lambda-a}
 \bigl(\log(1/\varepsilon)\bigr)^{\theta-1},
 &&\lambda<a, \label{face:eq:subcritical}\\
 I_a(\varepsilon)&\sim
 \frac{c_Z}{\theta\Gamma(a)}
 \bigl(\log(1/\varepsilon)\bigr)^\theta,
 &&\lambda=a, \label{face:eq:critical}\\
 I_a(\varepsilon)&\uparrow
 \int_Fh^{-a}\,d\sigma_F<\infty,
 &&\lambda>a. \label{face:eq:supercritical}
\end{align}
\end{lemma}

\begin{proof}
The gamma integral and Tonelli's theorem give
\begin{equation}
 \label{face:eq:gamma-transform}
 I_a(\varepsilon)=\frac1{\Gamma(a)}
 \int_0^\infty s^{a-1}e^{-\varepsilon s}Z_F(s)\,ds.
\end{equation}
Let $p=\theta-1$.  By the Laplace asymptotic and boundedness of $Z_F$
near zero, for every $s>0$,
$Z_F(s)\le C s^{-\lambda}(1+\log(2+s))^p$.
If $a>\lambda$, substitute $s=u/\varepsilon$ in
\eqref{face:eq:gamma-transform}.  For each $u>0$, the integrand divided
by $c_Z\varepsilon^{\lambda-a}(\log(1/\varepsilon))^p$
converges to $u^{a-\lambda-1}e^{-u}$.  For
$\varepsilon\le e^{-1}$ it is bounded by a constant times
$u^{a-\lambda-1}e^{-u}(2+\log(2+u))^p$.
This function is integrable.  Dominated convergence gives
\eqref{face:eq:subcritical}.

If $a=\lambda$, the part of the integral below a fixed $S>1$ is
bounded.  Above $S$, use
$Z_F(s)=c_Zs^{-a}(\log s)^p(1+o(1))$.  Writing
$L=\log(1/\varepsilon)$, one has
\[
 \int_S^\infty e^{-\varepsilon s}\frac{(\log s)^p}{s}\,ds
 =\frac{L^{p+1}}{p+1}+O(L^p+1).
\]
Indeed, replacing $e^{-\varepsilon s}$ by one on $[S,1/\varepsilon]$
changes the integral by at most
$\varepsilon\int_S^{1/\varepsilon}(\log s)^p\,ds=O(L^p)$.
On $[1/\varepsilon,\infty)$, the substitution
$u=\varepsilon s$ bounds the integral by $O(L^p+1)$.
Bounding the $o(1)$ remainder above a sufficiently large fixed $S$
by an arbitrary constant proves \eqref{face:eq:critical}.
Finally, if $a<\lambda$, the integrand in
\eqref{face:eq:gamma-transform} with $\varepsilon=0$ is integrable
both at zero and infinity.  Tonelli and monotone convergence give
\eqref{face:eq:supercritical}.
\end{proof}

For each positive-dimensional face, define a rate without assigning
an RLCT to the identically-zero case:
\begin{equation}
 \label{face:eq:face-rate}
 R_F(\varepsilon)=
 \begin{cases}
  \varepsilon^{-d_F/2},&m|_F\equiv0,\\
  1,&\min_Fm>0,\\
  \varepsilon^{\lambda_F-d_F/2}
   (\log(1/\varepsilon))^{\theta_F-1},
    &m|_F\not\equiv0,\ \min_Fm=0,\ \lambda_F<d_F/2,\\
  (\log(1/\varepsilon))^{\theta_F},
    &m|_F\not\equiv0,\ \min_Fm=0,\ \lambda_F=d_F/2,\\
  1,&m|_F\not\equiv0,\ \min_Fm=0,\ \lambda_F>d_F/2.
 \end{cases}
\end{equation}

\begin{corollary}[Facewise analytic rate]
 \label{face:cor:analytic-rate}
Under the assumptions of \cref{face:thm:characterization}, with $f$
real analytic near $D$, each certificate quantity and the dyadic node
count satisfy
\begin{equation}
 \label{face:eq:analytic-rate}
 N_{\mathrm{cover}}(\varepsilon)
 \asymp N_{\mathrm{rect}}(\varepsilon)
 \asymp N_{\mathrm{tree}}(\varepsilon)
 \asymp |T_{\mathrm{dyad}}(\varepsilon)|
 \asymp \max\left\{1,\max_{F\in\mathcal F_+(D)}R_F(\varepsilon)\right\}.
\end{equation}
\end{corollary}

\begin{proof}
Apply \cref{face:lem:analytic-integral} with $a=d_F/2$ whenever the
restriction is nonzero and has a zero.  If $m|_F\equiv0$, its integral
is exactly $\sigma_F(F)\varepsilon^{-d_F/2}$.  If $\min_Fm>0$, its
integral converges to a positive finite value.  There are finitely many
faces, so their sum is comparable to the largest term, with a baseline
one.  Now use \cref{face:thm:characterization}.
\end{proof}

The equality case $\lambda_F=d_F/2$ has one more logarithm than the
Laplace asymptotic.  The theorem retains this case even when
$\theta_F>1$; it needs no claim about whether such a face can dominate
all of its subfaces.  The constants in
\cref{face:lem:analytic-integral} are leading constants for the
integrals.  They do not assert that a discrete tree count has a limiting
leading constant.

\subsection{When the full-dimensional integral suffices}

\begin{proposition}[Full-box reduction]
 \label{face:prop:full-box}
Under the assumptions of \cref{face:thm:characterization}, either of the
following conditions implies
\begin{equation}
 \label{face:eq:full-box}
 N_{\mathrm{cover}}(\varepsilon)
 \asymp N_{\mathrm{rect}}(\varepsilon)
 \asymp N_{\mathrm{tree}}(\varepsilon)
 \asymp |T_{\mathrm{dyad}}(\varepsilon)|
 \asymp 1+I_D(\varepsilon):
\end{equation}
\begin{enumerate}[label=(\roman*)]
 \item Every minimizer lies in the interior of $D$.
 \item $m$ is nonnegative on a fixed open neighborhood of $D$, and its
 gradient is Lipschitz on a smaller closed neighborhood of $D$.
\end{enumerate}
\end{proposition}

\begin{proof}
In case (i), compactness gives $\min_{\partial D}m>0$, so every proper
face integral is bounded as $\varepsilon\downarrow0$.  Case (ii) is
proved in \cref{app:face:full-box}: nonnegativity beyond the box makes
low-sublevel points on each face centers of full-dimensional inward
cubes of side proportional to the square root of the sublevel height.
Packing these cubes and integrating their sublevel volumes gives
$I_F\le C(1+I_D)$ for every proper face.  In either case, use
\cref{face:thm:characterization}.
\end{proof}

\begin{corollary}[Interior analytic minimizers]
 \label{face:cor:interior}
Suppose $f$ is analytic near $D$ and every minimizer is interior.  Let
$(\lambda,\theta)$ be the RLCT pair on $D$.  Then $\lambda\le n/2$.
If $\lambda<n/2$, all the quantities in
\eqref{face:eq:full-box} have rate
$\varepsilon^{\lambda-n/2}(\log(1/\varepsilon))^{\theta-1}$.
Equality $\lambda=n/2$ holds exactly when every minimizer has positive
definite Hessian.  In that case the minimizers are finite,
$\theta=1$, and the rate is $\log(1/\varepsilon)$.  Moreover,
\begin{equation}
 \label{face:eq:morse-integral}
 I_D(\varepsilon)\sim
 \frac{(2\pi)^{n/2}}{\Gamma(n/2)}
 \left(\sum_{x\in\argmin_D f}\det(\nabla^2 f(x))^{-1/2}\right)
 \log(1/\varepsilon).
\end{equation}
\end{corollary}

\begin{proof}
At an interior minimizer $x_*$, the gradient vanishes and the Hessian is
positive semidefinite.  Taylor's theorem gives
$m(x_*+h)\le C\|h\|^2$, so the sublevel set of height $t$ contains a
ball of radius $c\sqrt t$.  Hence
$Z_D(s)\ge c s^{-n/2}$ and $\lambda\le n/2$.
If the Hessian has a unit null vector $v$, write $h=w+rv$ with $w\perp v$.
Analyticity gives, for small $h$,
\[
 m(x_*+w+rv)\le C(\|w\|^2+|r|^3).
\]
For example, this follows by bounding the quadratic term on $v^\perp$
and the Taylor remainder $O(\|h\|^3)$.
The corresponding anisotropic sublevel box has volume at least
$ct^{(n-1)/2+1/3}$.  Thus
$Z_D(s)\ge cs^{-((n-1)/2+1/3)}$, forcing
$\lambda\le(n-1)/2+1/3<n/2$.

If all Hessians are positive definite, each minimizer is isolated.
Their closed set is compact, and an infinite set would have a
nonisolated accumulation point, so there are finitely many.  Outside
fixed disjoint neighborhoods of these points, $m$ has a positive lower
bound.  At each minimizer, rescale $h=s^{-1/2}z$ in its Laplace integral.
The local quadratic lower bound gives an integrable Gaussian majorant,
and Taylor expansion and dominated convergence give
\[
 Z_D(s)\sim (2\pi)^{n/2}s^{-n/2}
 \sum_{x\in\argmin_D f}\det(\nabla^2f(x))^{-1/2}.
\]
Therefore $\lambda=n/2$ and $\theta=1$.
Apply \cref{face:lem:analytic-integral,face:prop:full-box} to obtain
the stated rate and \eqref{face:eq:morse-integral}.
\end{proof}

Under condition (ii) of \cref{face:prop:full-box}, the same
full-dimensional RLCT rate applies when $m$ is nonzero.  Every zero is
then a critical point in an open neighborhood.  A small ball centered
at that zero meets $D$ in a set containing a fixed fraction of its
volume, so the preceding quadratic upper bound again gives
$\lambda\le n/2$.  Boundary degeneracy need not satisfy the interior
Hessian equivalence just proved.

\subsection{Examples and the role of the boundary}

\begin{example}[An isolated degenerate minimum]
 \label{face:ex:quartic}
On $D=[-1,1]$, let $m(x)=x^4$ and fix a positive coefficient in
\eqref{face:eq:oracle}.  The minimum is isolated, but
\[
 I_D(\varepsilon)
 =\varepsilon^{-1/4}
 \int_{-\varepsilon^{-1/4}}^{\varepsilon^{-1/4}}
       (1+z^4)^{-1/2}\,dz
 \asymp\varepsilon^{-1/4}.
\]
The limiting integral over $\R$ is finite and positive.  Thus all three
certificate complexities have rate $\varepsilon^{-1/4}$, although the
optimal set has dimension zero.  Replacing $x^4$ by $x^2$ gives the
critical logarithmic rate.  A rule based only on the dimension of the
optimal set therefore requires a growth hypothesis, such as quadratic
growth; it does not apply to arbitrary isolated analytic minima.
\end{example}

\begin{example}[A singular intersection gives a logarithmic factor]
 \label{face:ex:product}
Let $m(x,y)=x^2y^2$ on $[-1,1]^2$.  This function is nonnegative on
$\R^2$, so \cref{face:prop:full-box} applies.  Writing $t=\delta^2$,
the sublevel area for $0<\delta<1$ is exactly
\[
 4\left(\int_0^\delta 1\,dx+
          \int_\delta^1\frac\delta x\,dx\right)
 =4\delta(1+\log(1/\delta)).
\]
It has rate $t^{1/2}\log(1/t)$.  Directly,
$\zeta_D(z)=4/(1-2z)^2$ for $\operatorname{Re}z<1/2$, so
$(\lambda,\theta)=(1/2,2)$.
The regularized integral can also be evaluated directly:
\[
 I_D(\varepsilon)
 =\frac4{\sqrt\varepsilon}
 \int_0^1\frac{\arctan(x/\sqrt\varepsilon)}x\,dx
 \sim\pi\varepsilon^{-1/2}\log(1/\varepsilon).
\]
To check the last asymptotic, split the integral at
$x=\sqrt\varepsilon$.  The lower part is bounded after rescaling.
On the upper part, replacing the arctangent by $\pi/2$ changes it by
at most $\sqrt\varepsilon\int_{\sqrt\varepsilon}^1x^{-2}\,dx\le1$.
Thus the certificate rate is
$\varepsilon^{-1/2}\log(1/\varepsilon)$.
\end{example}

\begin{example}[The full-box RLCT misses the dominant face]
 \label{face:ex:boundary}
On the four-dimensional cube
\[
 D=[0,0.9]\times[-0.4,0.5]^3,
 \qquad m(x,y,z,w)=x(1-x)+y^4+z^4+w^4,
\]
the unique minimizer is the origin, on the face $F=\{x=0\}\cap D$.
Taking $\alpha_i=1$ in \eqref{face:eq:oracle} makes every node
underestimator convex, since the Hessian of $m$ is
$\diag(-2,12y^2,12z^2,12w^2)$.

The Laplace integral factors.  For the first coordinate, substituting
$x=r/s$ and using $x(1-x)\ge x/10$ gives
\[
 \int_0^{0.9}e^{-s x(1-x)}\,dx\sim s^{-1}.
\]
For each quartic coordinate, substituting $y=s^{-1/4}v$ gives
\[
 \int_{-0.4}^{0.5}e^{-sy^4}\,dy
 \sim c_4s^{-1/4},\qquad
 c_4=\int_\R e^{-v^4}\,dv=2\Gamma(5/4).
\]
Consequently the full-box pair is $(\lambda_D,\theta_D)=(7/4,1)$,
and $I_D(\varepsilon)\asymp\varepsilon^{-1/4}$.
On $F$, the pair is $(3/4,1)$ and $d_F=3$, so
$I_F(\varepsilon)\asymp\varepsilon^{-3/4}$.
Every proper face other than $F$ has positive minimum: it fixes $x=0.9$
or at least one quartic coordinate at a nonzero endpoint.  The
face characterization therefore gives
\[
 N_{\mathrm{cover}}(\varepsilon)
 \asymp N_{\mathrm{rect}}(\varepsilon)
 \asymp N_{\mathrm{tree}}(\varepsilon)
 \asymp |T_{\mathrm{dyad}}(\varepsilon)|
 \asymp\varepsilon^{-3/4}.
\]
This also explains why nonnegativity merely on the root box is enough
for the facewise theorem but not for the full-box reduction:
$x(1-x)$ becomes negative immediately outside the face $x=0$.
\end{example}

The analytic input describes how sublevel measure concentrates near a
singular zero.  The additional optimization argument identifies which
intrinsic dimension must be paired with that concentration.  For this
oracle it is the maximum over root-box faces, and at boundary minima
that maximum can be attained on a proper face.  For face-exact
relaxations such as termwise McCormick bounds, the gap has a different
zero set; the geometric tools in the preceding section apply to that
separate setting.
